\chapter{NP-Completeness}\label{ch:npcomplete}

\begin{center}\itshape
Thousands of problems that look nothing alike are one problem in disguise:\\
solve any of them quickly, and you have solved them all.\\
\upshape(the content of this chapter, in our words)
\end{center}

On 20 March 1956 Kurt Gödel wrote a letter to John von Neumann, who was in hospital in Washington, dying of cancer. After wishing him a recovery, Gödel asked a question that nobody had asked before. Take a formula of first-order logic and a number $n$; deciding whether the formula has a proof of length at most $n$ is a finite problem, since one could try all proofs. Let $\varphi(n)$ be the number of steps the best machine needs for it, in the worst case. Gödel remarked that one can show $\varphi(n)\ge kn$ for some constant $k$, and asked how fast $\varphi(n)$ grows for the best machine. If $\varphi(n)$ grew like $kn$ or $kn^2$, he wrote (in our words), it would have consequences of the greatest importance: the thinking of a mathematician on yes-or-no questions could be entirely replaced by machines, despite the undecidability of the Entscheidungsproblem. He did not think this out of the question. Von Neumann, too ill, never answered, and died in February 1957. The letter became known to complexity theorists only when Juris Hartmanis published it, with a translation, in 1989; by then the question had a name.

The name came in May 1971, at the third Symposium on the Theory of Computing in Shaker Heights, Ohio. Stephen Cook, who had just moved to Toronto after the mathematics department of Berkeley declined to give him tenure (Richard Karp later called it a lasting shame for Berkeley), presented a paper on \emph{the complexity of theorem-proving procedures}. He proved that every problem solvable by a nondeterministic machine in polynomial time reduces, in polynomial time, to deciding whether a propositional formula is a tautology. If tautologies could be recognised quickly, so could the solutions of every problem whose solutions can be \emph{checked} quickly. In Moscow, Leonid Levin, a student of Kolmogorov, found the same phenomenon for six search problems, among them tiling; his two-page paper appeared in 1973, and he emigrated to the United States in 1978.

A year after Cook, at a symposium at IBM's research centre in Yorktown Heights, Richard Karp showed that twenty-one classical problems of combinatorics and operations research, which had resisted the efforts of many people for decades, all have this property: the travelling salesman, colouring, cliques, covers, knapsacks, integer programming. The theory of NP-completeness was born, and spread to every science that computes. Michael Garey and David Johnson's guide of 1979 listed some three hundred NP-complete problems; there are now thousands. Cook received the Turing Award in 1982, Karp in 1985. Whether the problems of this chapter can be solved in polynomial time, the question $\mathsf P=\mathsf{NP}$ that Gödel's letter anticipated, is one of the seven Millennium Prize Problems of the Clay Mathematics Institute, with a prize of a million dollars, offered in 2000 and still unclaimed.

\begin{think}
A Sudoku grid of size $n^2\times n^2$, half filled: how long does it take you to check a proposed solution, and how long to find one? Is it conceivable that every puzzle whose solutions are easy to check also has solutions that are easy to find?
\end{think}

\section{Certificates and reductions}\label{sec:certificates}

The class $\mathsf{NP}$ of Chapter~\ref{ch:timespace} was defined with nondeterministic machines. Its meaning is clearer with \emph{certificates}: a problem is in $\mathsf{NP}$ when every yes-instance has a short proof of being one, which can be checked quickly. A graph is $3$-colourable if a colouring is shown; a formula is satisfiable if an assignment is shown; a number is composite if a factor is shown.

\begin{proposition}\label{prop:certificates}
A language $L$ is in $\mathsf{NP}$ if and only if there are a polynomial $p$ and a deterministic machine $V$ (the \emph{verifier}) running in polynomial time such that, for every word $x$,
\[
x\in L\iff V\text{ accepts }\langle x,y\rangle\text{ for some }y\in\{0,1\}^{p(|x|)}.
\]
Such a $y$ is a \emph{certificate} for $x$.
\end{proposition}

\begin{proof}
Suppose that a nondeterministic machine $N$ decides $L$ in time $n^k$, and that each of its steps chooses among at most $2^b$ moves. A computation on $x$ is described by the sequence of its choices, written in binary with $b$ bits for each step: a word $y$ of length $p(|x|)=b|x|^k$. The verifier $V$ simulates $N$ on $x$, making at each step the choice described by the next $b$ bits of $y$ (and rejecting if they name no possible move), and accepts if this computation accepts. It runs in polynomial time, and some $y$ makes it accept if and only if some computation of $N$ accepts $x$.

Conversely, given $V$ and $p$, a nondeterministic machine on $x$ writes a word $y$ of length $p(|x|)$, choosing each bit, and runs $V$ on $\langle x,y\rangle$. It runs in polynomial time and accepts $x$ if and only if some certificate exists.
\end{proof}

Every problem in $\mathsf P$ is in $\mathsf{NP}$ (ignore the certificate), and $\mathsf{NP}\subseteq\mathsf{EXP}$ (try all $2^{p(n)}$ certificates). Problems in $\mathsf{coNP}$ have short certificates for \emph{no}-instances: a formula is a tautology if no assignment falsifies it, and nobody knows short certificates for that.

To compare problems we use the reductions of Chapter~\ref{ch:turing}, restricted to polynomial time.

\begin{definition}\label{def:polyreduction}
A language $A$ \emph{reduces in polynomial time} to $B$, written $A\le_pB$, if there is a function $f$ computable in polynomial time with $x\in A\iff f(x)\in B$ for every $x$. A language $B$ is \emph{NP-hard} if $A\le_pB$ for every $A\in\mathsf{NP}$, and \emph{NP-complete} if moreover $B\in\mathsf{NP}$.
\end{definition}

An optimisation problem, such as finding a largest clique, is called NP-hard when its decision version (is there a clique of size at least $k$?) is: an algorithm for the optimisation problem answers the decision question at once.

\begin{proposition}\label{prop:reductions}
\begin{enumerate}
\item If $A\le_pB$ and $B\le_pC$, then $A\le_pC$.
\item If $A\le_pB$ and $B\in\mathsf P$, then $A\in\mathsf P$.
\item If $B$ is NP-complete, then $\mathsf P=\mathsf{NP}$ if and only if $B\in\mathsf P$.
\item If $B$ is NP-complete, $C\in\mathsf{NP}$ and $B\le_pC$, then $C$ is NP-complete.
\end{enumerate}
\end{proposition}

\begin{proof}
(1) If $f$ runs in time $O(n^a)$ and $g$ in time $O(n^b)$, then $f(x)$ has length $O(|x|^a)$, and $g(f(x))$ is computed in time $O(|x|^a+|x|^{ab})$. (2) To decide $A$, compute $f(x)$ and run the algorithm for $B$ on it: polynomial time, for the same reason. (3) If $B\in\mathsf P$, every $A\in\mathsf{NP}$ reduces to $B$ and is in $\mathsf P$ by (2); conversely $B\in\mathsf{NP}=\mathsf P$. (4) Every $A\in\mathsf{NP}$ reduces to $B$, hence to $C$ by (1).
\end{proof}

Point (4) is how NP-completeness spreads: once one problem is known to be NP-complete, a single reduction from it proves the next one. The hard part is the first.

\section{The theorem of Cook and Levin}\label{sec:cooklevin}

A \emph{formula in conjunctive normal form} (CNF) is a conjunction of \emph{clauses}, each a disjunction of literals, as in Tome~C1, Chapter~1. \textup{\textsc{Sat}} is the set of satisfiable CNF formulas, written in some reasonable way as words, and \textup{\textsc{3-Sat}} the set of satisfiable formulas whose clauses have at most three literals. Both are in $\mathsf{NP}$: a satisfying assignment is a certificate, checked clause by clause.

\begin{theorem}[Cook 1971, Levin 1973]\label{thm:cooklevin}
\textup{\textsc{Sat}} is NP-complete.
\end{theorem}

\begin{proof}
Let $L\in\mathsf{NP}$, with a verifier $V$ and a polynomial $p$ as in Proposition~\ref{prop:certificates}. By Theorem~\ref{thm:multitape} we may take $V$ with one tape; on the input $x\#y$, with $|x|=n$ and $|y|=p(n)$, it halts within $T=T(n)$ steps, a polynomial in $n$. We also let $V$ stay where it is after halting: for the halting states $h\in\{q_{\mathrm{acc}},q_{\mathrm{rej}}\}$ we set $\Delta(h,s)=(h,s,0)$, and for the other states $\Delta(q,s)=(q',s',D)$ with $D=-1$ for a left move and $D=+1$ for a right move, as given by $\delta$. Then the computation has a configuration $C_t$ for every $t=0,1,\dots,T$, and $V$ accepts if and only if $C_T$ is in state $q_{\mathrm{acc}}$. In $T$ steps the head stays in the cells $-T,\dots,T$, and the other cells keep their blanks.

Given $x$, we build a CNF formula $\varphi_x$ that describes the table of all the configurations $C_0,\dots,C_T$ (Figure~\ref{fig:tableau}). Its variables, for $0\le t\le T$, $-T\le i\le T$, $s\in\Gamma$ and $q\in Q$, are
\begin{itemize}
\item $X_{t,i,s}$: at time $t$, cell $i$ contains $s$;
\item $H_{t,i}$: at time $t$, the head is on cell $i$;
\item $S_{t,q}$: at time $t$, the state is $q$.
\end{itemize}
The formula $\varphi_x$ is the conjunction of the following clauses.
\begin{enumerate}
\item \emph{Exactly one.} For each $t$ and $i$, the clause $\bigvee_sX_{t,i,s}$ and the clauses $\neg X_{t,i,s}\lor\neg X_{t,i,s'}$ for $s\ne s'$; in the same way exactly one $H_{t,i}$ is true for each $t$, and exactly one $S_{t,q}$.
\item \emph{Start.} The unit clauses $S_{0,q_0}$ and $H_{0,0}$; $X_{0,i,x_{i+1}}$ for $0\le i<n$; $X_{0,n,\#}$; for $n<i\le n+p(n)$, the unit clauses $\neg X_{0,i,s}$ for every $s\notin\{0,1\}$; and $X_{0,i,\sqcup}$ for the other cells.
\item \emph{Moves.} For each $t<T$, each cell $i$, each $q\in Q$ and $s\in\Gamma$, with $\Delta(q,s)=(q',s',D)$: if $-T\le i+D\le T$, the three clauses
\[
\bigl(H_{t,i}\land S_{t,q}\land X_{t,i,s}\bigr)\to S_{t+1,q'},\qquad
\bigl(H_{t,i}\land S_{t,q}\land X_{t,i,s}\bigr)\to X_{t+1,i,s'},
\]
\[
\bigl(H_{t,i}\land S_{t,q}\land X_{t,i,s}\bigr)\to H_{t+1,i+D},
\]
each of which is the clause $\neg H_{t,i}\lor\neg S_{t,q}\lor\neg X_{t,i,s}\lor\ell$; otherwise the clause $\neg H_{t,i}\lor\neg S_{t,q}\lor\neg X_{t,i,s}$.
\item \emph{Frame.} For each $t<T$, $i$ and $s$, the clause $H_{t,i}\lor\neg X_{t,i,s}\lor X_{t+1,i,s}$: a cell away from the head keeps its symbol.
\item \emph{Accept.} The unit clause $S_{T,q_{\mathrm{acc}}}$.
\end{enumerate}
The cells $n+1,\dots,n+p(n)$ of the first row are left free, apart from being $0$ or $1$: they hold the certificate.

\emph{Size.} There are $O(T^2)$ variables. The exactly-one clauses for the head take $O(T^2)$ clauses for each $t$, and every other group has $O(1)$ clauses for each triple $(t,i,\cdot)$, the alphabet and the states being fixed; so $\varphi_x$ has $O(T^3)$ clauses of length $O(T)$, and a machine writes it in time polynomial in $n$, since $T$ and $p(n)$ are polynomials and each clause is written by a simple loop.

\emph{If $x\in L$}, take a certificate $y$, run $V$ on $x\#y$, and give each variable the truth value it has in this computation. Every clause holds, by the definition of a step; the last because $V$ accepts.

\emph{If $\varphi_x$ is satisfied by some assignment}, the exactly-one clauses give, for each $t$, a state $q_t$, a head position $i_t$, and a word on the cells $-T,\dots,T$: a configuration $D_t$ (with blanks outside). Let $y$ be the word written in the cells $n+1,\dots,n+p(n)$ of $D_0$; it is in $\{0,1\}^{p(n)}$, and $D_0$ is the initial configuration of $V$ on $x\#y$. Suppose that $D_t$ is the configuration $C_t$ of this computation, and let $s$ be the symbol under the head and $\Delta(q_t,s)=(q',s',D)$. The move clauses for $(t,i_t,q_t,s)$ force $q_{t+1}=q'$, cell $i_t$ to contain $s'$ at time $t+1$, and $i_{t+1}=i_t+D$ (if $i_t+D$ were outside $[-T,T]$, the clause $\neg H_{t,i_t}\lor\neg S_{t,q_t}\lor\neg X_{t,i_t,s}$ of $\varphi_x$ would be false). The frame clauses keep the other cells. So $D_{t+1}=C_{t+1}$. By induction $D_T=C_T$, which is in state $q_{\mathrm{acc}}$ by the last clause: $V$ accepts $x\#y$, and $x\in L$.

So $x\mapsto\varphi_x$ is a polynomial-time reduction of $L$ to \textup{\textsc{Sat}}.
\end{proof}

\begin{figure}[ht]
\centering
\fig{ch09-tableau}
\caption{The table of the proof of Theorem~\ref{thm:cooklevin}: one row per time step, one column per cell, the head shaded in each row. The first row holds the input $x$, a separator, and the certificate $y$, whose cells are left free. The clauses say that each row follows from the previous one by a step of the verifier, a statement about a few neighbouring cells at a time, and that the last row is accepting. A satisfying assignment exists exactly when some certificate is accepted.}
\label{fig:tableau}
\end{figure}

The proof shows more than the statement. The formula is built by a simple program from $V$ and $x$, and its satisfying assignments correspond to accepted certificates: given one, the certificate is read off the first row. Cook's original statement was about tautologies (the complement of satisfiability, a formula being a tautology if its negation is unsatisfiable), hence the title \emph{theorem-proving procedures}.

\begin{theorem}[Karp 1972]\label{thm:threesat}
\textup{\textsc{3-Sat}} is NP-complete.
\end{theorem}

\begin{proof}
We reduce \textup{\textsc{Sat}} to it. A clause $\ell_1\lor\dots\lor\ell_k$ with $k\ge4$ is replaced, with new variables $z_1,\dots,z_{k-3}$ used for this clause only, by the clauses
\[
(\ell_1\lor\ell_2\lor z_1)\land(\neg z_1\lor\ell_3\lor z_2)\land\dots\land(\neg z_{k-3}\lor\ell_{k-1}\lor\ell_k).
\]
If an assignment makes some $\ell_j$ true, it extends: set $z_i$ true for the clauses before the one containing $\ell_j$, and false after. If an assignment of the new formula made all the $\ell_j$ false, then $z_1$ would be true by the first clause, then $z_2$ by the second, and so on, and the last clause would be false. So each long clause can be replaced while preserving satisfiability, and the new formula has size linear in the old one.
\end{proof}

We may further simplify: a clause containing $x$ and $\neg x$ is always true and can be dropped, and repeated literals in a clause can be merged. So we assume below that \emph{the literals of each clause involve distinct variables}; then each clause has one, two or three literals.

\section{A web of reductions}\label{sec:karp}

Figure~\ref{fig:reductions} shows the reductions of this section. Each takes a formula $\varphi$ in 3-CNF, with variables $x_1,\dots,x_n$ and clauses $c_1,\dots,c_m$, and builds an instance out of \emph{gadgets}: pieces that imitate a variable, which must take one of two values, and pieces that imitate a clause, which must be satisfied by one of its literals. Every problem of the figure is in $\mathsf{NP}$, with the obvious certificate (a set of vertices, a colouring, a cycle, a subset, a vector), except that for integer programming this needs an argument.

\begin{figure}[ht]
\centering
\fig{ch09-reductions}
\caption{The reductions of this chapter. An arrow from $A$ to $B$ is a polynomial reduction $A\le_pB$, so that every problem below \textup{\textsc{Sat}} is NP-hard. Karp's paper of 1972 had a similar tree, with twenty-one problems.}
\label{fig:reductions}
\end{figure}

\subsection*{Independent sets, cliques and covers}

An \emph{independent set} in a graph is a set of pairwise non-adjacent vertices, a \emph{clique} a set of pairwise adjacent vertices, and a \emph{vertex cover} a set of vertices meeting every edge. The problems \textup{\textsc{Independent set}}, \textup{\textsc{Clique}} and \textup{\textsc{Vertex cover}} ask, for a graph $G$ and a number $k$, whether $G$ has an independent set or a clique of size at least $k$, or a vertex cover of size at most $k$.

\begin{theorem}\label{thm:isclique}
\textup{\textsc{Independent set}}, \textup{\textsc{Clique}} and \textup{\textsc{Vertex cover}} are NP-complete.
\end{theorem}

\begin{proof}
Given $\varphi$, build the graph $G_\varphi$ with one vertex for each occurrence of a literal in a clause; join the vertices of each clause by edges (a triangle, an edge or nothing), and join any two vertices labelled by complementary literals $x_i$ and $\neg x_i$. We claim that $\varphi$ is satisfiable if and only if $G_\varphi$ has an independent set of size $m$.

If an assignment satisfies $\varphi$, choose in each clause one true literal. The $m$ chosen vertices lie in different clauses, and are never complementary since they are all true: they form an independent set. Conversely an independent set of size $m$ contains at most one vertex per clause, hence exactly one, and never two complementary literals. Make the literals of these vertices true (which is consistent), and the other variables arbitrary: every clause has a true literal.

A set $S$ is independent in $G$ if and only if it is a clique in the complement graph $\bar G$, and if and only if $V\setminus S$ is a vertex cover of $G$. So $(G,k)\mapsto(\bar G,k)$ reduces \textup{\textsc{Independent set}} to \textup{\textsc{Clique}}, and $(G,k)\mapsto(G,|V|-k)$ reduces it to \textup{\textsc{Vertex cover}}.
\end{proof}

So the vertex cover problem, for which Tome~C2, Chapter~12, gives a $2$-approximation, is NP-hard. In the language of Tome~C1, Chapter~8: computing the clique number $\omega(G)$ or the independence number $\alpha(G)$ is NP-hard, and Ramsey's theorem says that one of them is always at least logarithmic.

\subsection*{Colourings}

\begin{theorem}\label{thm:threecol}
Deciding whether a graph is $3$-colourable is NP-complete.
\end{theorem}

\begin{proof}
Build a graph from $\varphi$ as follows (Figure~\ref{fig:colouring}). A triangle $T,F,B$, the \emph{palette}, gets three different colours, which we call true, false and base. For each variable $x_i$, two vertices $x_i$ and $\neg x_i$, joined to each other and to $B$: in a colouring they receive the colours of $T$ and $F$ in some order, which defines a truth assignment. For a clause $a\lor b\lor c$ (a clause with fewer literals repeats one), add the gadget of the figure: a triangle $p,q,o$ with $p$ joined to $a$ and $q$ to $b$; a triangle $p',q',o'$ with $p'$ joined to $o$ and $q'$ to $c$; and $o'$ joined to $F$ and $B$.

The key is the triangle $p,q,o$ with inputs $u$ (joined to $p$) and $v$ (joined to $q$), whose colours are true or false. If both $u$ and $v$ are false, then $p$ and $q$ avoid false, so they take true and base, and $o$ must be false. If $u$ is true, $o$ can be true: colour $p$ false and $q$ base. If $v$ is true, $o$ can be true: colour $q$ false and $p$ base. So the output $o$ can be coloured true if and only if an input is true, and is forced to be false otherwise.

If an assignment satisfies $\varphi$, colour the palette and the literals accordingly. In a clause $a\lor b\lor c$, if $a$ or $b$ is true, colour $o$ true, then $o'$ true (its input $o$ is true); if $a$ and $b$ are false, $o$ is forced false, and $c$ is true, so $o'$ can again be true. Since $o'$ is true, it differs from $F$ and $B$, and the colouring is proper. Conversely, in a proper $3$-colouring, $o'$ is adjacent to $F$ and $B$ and so has the colour of $T$, true. If $a$, $b$ and $c$ were all false, $o$ would be forced false, and then $o'$ too, a contradiction. So the assignment read on the literals satisfies every clause. The graph has $3+2n+6m$ vertices.
\end{proof}

\begin{figure}[ht]
\centering
\fig{ch09-colouring}
\caption{Gadgets for $3$-colouring. Left: the palette triangle and a variable, whose two vertices receive the colours of $T$ and $F$ in some order. Right: the gadget of a clause $a\lor b\lor c$; the vertices $F$ and $B$ are those of the palette. Its output $o'$ can be coloured like $T$ if and only if one of $a$, $b$, $c$ is.}
\label{fig:colouring}
\end{figure}

This is the result announced in Tome~C1, Chapter~7. Deciding $2$-colourability is easy (a graph is bipartite if and only if it has no odd cycle), and $4$-colourability is NP-complete too (Exercise~\ref{ex:kcol}). Larry Stockmeyer showed in 1973 that $3$-colourability stays NP-complete for planar graphs (quoted), whereas every planar graph is $4$-colourable: for planar graphs the question ``$\chi\le4$?'' has the answer yes, and ``$\chi\le3$?'' is as hard as anything in $\mathsf{NP}$.

\subsection*{Hamiltonian cycles and the travelling salesman}

A \emph{Hamiltonian cycle} passes through every vertex exactly once. The problems \textup{\textsc{Directed Ham}} and \textup{\textsc{Ham cycle}} ask whether a directed or an undirected graph has one.

\begin{theorem}\label{thm:dham}
\textup{\textsc{Directed Ham}} is NP-complete.
\end{theorem}

\begin{proof}
For each variable $x_i$ build a row of $3m+3$ vertices $w_{i,0},\dots,w_{i,3m+2}$, with arrows in both directions between consecutive ones, and two vertices $s_i$ and $t_i$ with arrows $s_i\to w_{i,0}$, $s_i\to w_{i,3m+2}$, $w_{i,0}\to t_i$, $w_{i,3m+2}\to t_i$ (Figure~\ref{fig:hamrow}). Join the rows in a ring: $t_i\to s_{i+1}$, and $t_n\to s_1$. In row $i$, the pair $(w_{i,3j-1},w_{i,3j})$ belongs to clause $c_j$, and the vertices $w_{i,3j+1}$ and $w_{i,1}$ separate the pairs. For each clause $c_j$ add a vertex $c_j$; if $x_i$ appears in $c_j$, add the arrows $w_{i,3j-1}\to c_j\to w_{i,3j}$; if $\neg x_i$ appears, add $w_{i,3j}\to c_j\to w_{i,3j-1}$.

\emph{From an assignment to a cycle.} Go round the ring; in row $i$, go from $s_i$ along the row from left to right if $x_i$ is true, from right to left if it is false, and on to $t_i$. This visits every vertex except the $c_j$. For each clause $c_j$ choose a true literal. If it is $x_i$, row $i$ is traversed from left to right and uses the arrow $w_{i,3j-1}\to w_{i,3j}$; replace it by the detour $w_{i,3j-1}\to c_j\to w_{i,3j}$. If it is $\neg x_i$, row $i$ is traversed from right to left and uses $w_{i,3j}\to w_{i,3j-1}$, which the detour $w_{i,3j}\to c_j\to w_{i,3j-1}$ replaces. We get a Hamiltonian cycle.

\emph{From a cycle to an assignment.} Let $H$ be a Hamiltonian cycle. We first show that if $H$ enters $c_j$ from a vertex $a$ of row $i$, it leaves to the other vertex $b$ of the same pair. Take the case $a=w_{i,3j-1}$, $b=w_{i,3j}$ ($x_i$ in $c_j$; the other case is the mirror image), and suppose that $H$ leaves $c_j$ to some other vertex. The arrows into $b$ come from $a$, $c_j$ and $w_{i,3j+1}$; since $a$ and $c_j$ continue elsewhere, $H$ enters $b$ from $w_{i,3j+1}$, and then leaves it to $a$, its only other out-neighbour. So $H$ contains $b\to a\to c_j$. Now the separator $w_{i,3j-2}$ has arrows only to and from its two neighbours on the row, $w_{i,3j-3}$ and $a$; but $a$ already has its predecessor $b$ and its successor $c_j$, so $H$ must both enter and leave $w_{i,3j-2}$ through $w_{i,3j-3}$, which is impossible. Hence every visit to a clause vertex is a detour $a\to c_j\to b$ within one pair.

Replace each detour by the arrow $a\to b$: we get a Hamiltonian cycle $H'$ of the graph without clause vertices. In that graph, the vertices of row $i$ can only be entered from $s_i$, at one of the two ends, and left towards $t_i$, also from an end. A path that enters the row at one end and visits all its vertices before leaving must run straight to the other end. So $H'$ traverses each row from one end to the other; set $x_i$ true if it goes from left to right. A detour through $c_j$ from row $i$ used an arrow of the row in the direction $H'$ follows: from left to right if $x_i$ is in $c_j$, from right to left if $\neg x_i$ is. In both cases that literal is true, and $c_j$ is satisfied.

The graph has $n(3m+5)+m$ vertices and is built in polynomial time.
\end{proof}

\begin{figure}[ht]
\centering
\fig{ch09-hamrow}
\caption{The gadget of a variable $x_i$ in the proof of Theorem~\ref{thm:dham}, with $m=2$ clauses: a row that can be traversed in either direction (left to right means $x_i$ true), separators (shaded) between the pairs of the clauses, and detours through the clause vertices, which can only be taken in the direction matching the literal.}
\label{fig:hamrow}
\end{figure}

\begin{corollary}\label{cor:ham}
\textup{\textsc{Ham cycle}} is NP-complete, and so is the travelling salesman problem: given distances $d(u,v)\in\{1,2\}$ between $n$ cities and a bound $K$, is there a tour of length at most $K$?
\end{corollary}

\begin{proof}
Given a directed graph $D$, replace each vertex $v$ by a path of three vertices $v^{\mathrm{in}}$, $v^{\mathrm{mid}}$, $v^{\mathrm{out}}$, and each arrow $u\to v$ by an edge $u^{\mathrm{out}}v^{\mathrm{in}}$. A directed Hamiltonian cycle of $D$ gives an undirected one. Conversely, in an undirected Hamiltonian cycle each $v^{\mathrm{mid}}$, of degree~$2$, is passed through between $v^{\mathrm{in}}$ and $v^{\mathrm{out}}$; orient the cycle so that it goes $v^{\mathrm{in}}\to v^{\mathrm{mid}}\to v^{\mathrm{out}}$ for one $v$. The next vertex is some $u^{\mathrm{in}}$, since $v^{\mathrm{out}}$ has no other neighbours, so the cycle goes through $u$ in the same direction, and so on: it comes from a directed Hamiltonian cycle of $D$.

For the travelling salesman, given a graph on $n$ vertices, set $d(u,v)=1$ for edges and $2$ for non-edges, and $K=n$. A tour of length $n$ uses only edges, and is a Hamiltonian cycle.
\end{proof}

This is the reduction behind Tome~C2, Proposition~12.5, where the non-edges get a huge length instead of $2$ to show that the general travelling salesman problem cannot even be approximated. It also proves the claim of Tome~C1, Chapter~5, that no simple criterion for Hamiltonicity is to be expected: a criterion checkable in polynomial time would give $\mathsf P=\mathsf{NP}$.

\subsection*{Numbers: subset sum, knapsack, integer programming}

\textup{\textsc{Subset sum}} asks, for positive integers $a_1,\dots,a_N$ and a target $t$, written in binary, whether some subset of the $a_k$ adds up to $t$.

\begin{theorem}\label{thm:subsetsum}
\textup{\textsc{Subset sum}} is NP-complete.
\end{theorem}

\begin{proof}
Given $\varphi$, write numbers in base $10$ with $n+m$ digits, one column for each variable and one for each clause. For each variable $x_i$ take two numbers: $u_i$, with a $1$ in column $x_i$ and a $1$ in the column of each clause containing $x_i$; and $v_i$, with a $1$ in column $x_i$ and a $1$ in the column of each clause containing $\neg x_i$. For each clause $c_j$ take two \emph{slack} numbers $g_j=h_j$ with a single $1$, in column $c_j$. The target has digit $1$ in every variable column and $3$ in every clause column:
\[
t=\underbrace{1\cdots1}_{n}\underbrace{3\cdots3}_{m}.
\]
In each column the digits of all the numbers add up to at most $5$ (three literals and two slacks), so adding any subset produces no carries, and a subset sums to $t$ if and only if each column does.

If an assignment satisfies $\varphi$, take $u_i$ when $x_i$ is true and $v_i$ when it is false. Variable columns sum to $1$, and the column of $c_j$ sums to the number of true literals of $c_j$, between $1$ and $3$; add $2$, $1$ or $0$ slack numbers to reach $3$. Conversely, if a subset sums to $t$, each variable column forces exactly one of $u_i$, $v_i$; set $x_i$ true if $u_i$ is chosen. In the column of $c_j$ the slacks give at most $2$, so at least one chosen number comes from a literal of $c_j$, and that literal is true. The numbers have $n+m$ digits, so their binary length is polynomial.
\end{proof}

\begin{corollary}\label{cor:knapsack}
The knapsack problem is NP-complete: given items with weights $w_k$ and values $v_k$, a capacity $W$ and a goal $K$, is there a set of items of total weight at most $W$ and total value at least $K$?
\end{corollary}

\begin{proof}
Given $a_1,\dots,a_N$ and $t$, take items with $w_k=v_k=a_k$, and $W=K=t$. A set of weight at most $t$ and value at least $t$ has sum exactly $t$.
\end{proof}

So the dynamic programme of Tome~C2, Theorem~5.3, which runs in time $O(NW)$, is probably as good as one can hope for in a sense made precise in the next section: its time is polynomial in the \emph{value} of $W$, not in its length.

An \emph{integer program} asks, for an integer matrix $A$ and an integer vector $b$, whether some $x\in\Z^n$ satisfies $Ax\le b$.

\begin{theorem}\label{thm:ip}
Integer programming is NP-hard, even when the solutions are required to be in $\{0,1\}^n$.
\end{theorem}

\begin{proof}
Given $\varphi$, use one integer variable $x_i$ for each variable, with the constraints $0\le x_i\le1$. A clause becomes a linear inequality by writing $x_i$ for the literal $x_i$ and $1-x_i$ for $\neg x_i$, and requiring a sum of at least $1$: for example $x_1\lor\neg x_2\lor x_3$ becomes $x_1+(1-x_2)+x_3\ge1$. The solutions in $\Z^n$ are exactly the satisfying assignments.
\end{proof}

This keeps the promise of Tome~C2, Chapter~7: linear programming is in $\mathsf P$, but asking for integer points makes it NP-hard. Integer programming is in fact NP-complete, but membership in $\mathsf{NP}$ needs work, since a solution might a priori have huge entries: Christos Papadimitriou proved in 1981 that if a solution exists, one exists with entries of polynomially bounded length (quoted).

Many more problems join the club through similar gadgets. Computing the edge expansion $h(G)$ of Tome~C1, Chapter~9, is NP-hard (Bojan Mohar, 1989, quoted), which is why Cheeger's inequality, Tome~C1, Theorem~9.9, which approximates it by an eigenvalue, is so valuable. Exercises~\ref{ex:hampath} to~\ref{ex:longestpath} give a few more.

\section{Numbers in binary and in unary}\label{sec:strong}

Subset sum is NP-complete, yet the dynamic programme of Tome~C2 decides it in time $O(Nt)$. There is no contradiction: the input has length about $N\log t$, and $O(Nt)$ is exponential in $\log t$. The reduction of Theorem~\ref{thm:subsetsum} produces numbers with $n+m$ decimal digits, exponentially large in the size of $\varphi$. If the numbers are small, the problem is easy.

\begin{definition}\label{def:strong}
For an instance $I$ of a problem with numbers, let $|I|$ be its length and $\max(I)$ the largest number in it. An algorithm is \emph{pseudo-polynomial} if it runs in time polynomial in $|I|$ and $\max(I)$. A problem is \emph{strongly NP-hard} if, for some polynomial $q$, it remains NP-hard when restricted to the instances with $\max(I)\le q(|I|)$; equivalently, when the numbers are written in unary.
\end{definition}

\begin{proposition}\label{prop:strongpseudo}
If a strongly NP-hard problem has a pseudo-polynomial algorithm, then $\mathsf P=\mathsf{NP}$.
\end{proposition}

\begin{proof}
On the instances with $\max(I)\le q(|I|)$, the pseudo-polynomial algorithm runs in time polynomial in $|I|$ and $q(|I|)$, hence polynomial in $|I|$. So an NP-hard problem is in $\mathsf P$, and $\mathsf P=\mathsf{NP}$ by Proposition~\ref{prop:reductions}.
\end{proof}

So the problems split. Knapsack and subset sum have pseudo-polynomial algorithms, and are not strongly NP-hard unless $\mathsf P=\mathsf{NP}$: their hardness lives in the size of the numbers. The travelling salesman problem is strongly NP-hard: Corollary~\ref{cor:ham} uses only the distances $1$ and $2$. So are all the graph problems of Section~\ref{sec:karp}, which have no numbers to speak of apart from $k\le|V|$. Garey and Johnson proved in 1975 that \textup{\textsc{3-Partition}}, which asks whether $3k$ numbers can be split into $k$ triples with equal sums, is strongly NP-complete (quoted); it is the usual starting point for strong hardness of scheduling and packing problems.

The distinction decides which problems have the best kind of approximation, a \emph{fully polynomial approximation scheme} (FPTAS): an algorithm which, given $\varepsilon>0$, finds a solution within a factor $1+\varepsilon$ of the optimum in time polynomial in $|I|$ and $1/\varepsilon$. Knapsack has one (Tome~C2, Theorem~12.8).

\begin{proposition}\label{prop:nofptas}
Let $\Pi$ be an optimisation problem whose values are positive integers, with optimum $\mathrm{OPT}(I)\le r(|I|,\max(I))$ for a polynomial $r$. If $\Pi$ is strongly NP-hard and has an FPTAS, then $\mathsf P=\mathsf{NP}$.
\end{proposition}

\begin{proof}
Take a minimisation problem (maximisation is similar). Run the FPTAS with $\varepsilon=1/(r(|I|,\max(I))+1)$. It returns a value $A\le(1+\varepsilon)\mathrm{OPT}<\mathrm{OPT}+1$, since $\varepsilon\,\mathrm{OPT}<1$, and $A$ is an integer at least $\mathrm{OPT}$: so $A=\mathrm{OPT}$. The running time is polynomial in $|I|$ and $1/\varepsilon$, hence in $|I|$ and $\max(I)$: a pseudo-polynomial exact algorithm, and Proposition~\ref{prop:strongpseudo} applies.
\end{proof}

This is the remark of Tome~C2, Chapter~12, that strongly NP-hard problems such as the travelling salesman (even with the triangle inequality, where the distances $1$ and $2$ of Corollary~\ref{cor:ham} satisfy it) have no FPTAS unless $\mathsf P=\mathsf{NP}$. Whether they have approximation algorithms with \emph{some} constant ratio is another story, the story of the PCP theorem, told in Chapter~11.

\section{What NP-completeness means}\label{sec:meaning}

An NP-completeness proof is a statement of the form ``if this problem has a polynomial algorithm, then so does every problem in $\mathsf{NP}$''. Almost everyone believes that $\mathsf P\ne\mathsf{NP}$: in fifty years, thousands of NP-complete problems have been attacked by thousands of people, and a fast algorithm for any one of them would solve all the others. Most believe more: that $\mathsf{NP}\ne\mathsf{coNP}$, so that unsatisfiable formulas have no short certificates (Exercise~\ref{ex:conpcomplete}). Chapter~\ref{ch:timespace} explains why diagonalisation alone cannot settle the question, and Chapter~12 why the known methods of proving lower bounds cannot either.

Hardness is a statement about the worst case and about exact answers. In practice, solvers for \textup{\textsc{Sat}} handle formulas with millions of variables arising from hardware verification, and integer programming solvers schedule airlines; NP-hardness says only that no algorithm will handle every instance quickly. It also says nothing about cryptography, which needs problems that are hard \emph{on average}, a stronger and different requirement, the subject of Tome~C4.

Decision problems are enough, because finding a solution is no harder than deciding whether one exists.

\begin{proposition}\label{prop:selfreduce}
If \textup{\textsc{Sat}} is in $\mathsf P$, there is a polynomial algorithm which, given a satisfiable CNF formula, finds a satisfying assignment.
\end{proposition}

\begin{proof}
Given a satisfiable $\varphi(x_1,\dots,x_n)$, ask whether $\varphi$ with $x_1$ replaced by true (clauses containing $x_1$ deleted, literals $\neg x_1$ removed) is satisfiable. If yes, fix $x_1$ true and continue with that formula; if no, then $\varphi$ with $x_1$ false is satisfiable, since $\varphi$ is; fix $x_1$ false. After $n$ questions all variables are fixed, and each formula along the way stayed satisfiable, so the final assignment satisfies $\varphi$.
\end{proof}

Combined with Theorem~\ref{thm:cooklevin}, whose proof reads a certificate off a satisfying assignment, this shows that if $\mathsf P=\mathsf{NP}$ then for every problem in $\mathsf{NP}$ certificates can be \emph{found} in polynomial time: proofs of theorems of bounded length, as Gödel asked, factorisations of numbers, keys of ciphers.

Are there problems in $\mathsf{NP}$ that are neither in $\mathsf P$ nor NP-complete? If $\mathsf P=\mathsf{NP}$, no. Otherwise yes:

\begin{theorem}[Ladner 1975]\label{thm:ladner}
If $\mathsf P\ne\mathsf{NP}$, there is a language in $\mathsf{NP}\setminus\mathsf P$ that is not NP-complete. \emph{(Quoted.)}
\end{theorem}

Ladner's language is artificial: \textup{\textsc{Sat}}, with holes punched in it by a slow diagonalisation. Two natural problems are suspected of living in between: factoring integers, which is in $\mathsf{NP}\cap\mathsf{coNP}$ (a factorisation into primes is a certificate either way, primes being recognisable in polynomial time, Tome~C2, Chapter~10), so that it is not NP-complete unless $\mathsf{NP}=\mathsf{coNP}$; and graph isomorphism, for which László Babai announced in 2015 an algorithm in time $2^{(\log n)^{O(1)}}$ (quoted).

\begin{beauty}
Before 1971 the travelling salesman, the colouring of maps, the packing of knapsacks, the scheduling of jobs and the proving of theorems were separate problems, each with its specialists and its folklore of failed attempts. The Cook--Levin theorem showed that they are one problem. The proof is almost naive: a computation is a table, and ``this table is correct'' is a conjunction of facts about a few neighbouring cells, which is a formula. Through that remark a single open question, $\mathsf P=\mathsf{NP}$, took the place of thousands, and every new NP-completeness proof since, like the gadgets of this chapter, is a small piece of engineering showing how much one problem can imitate another.
\end{beauty}

\section{Exercises}

\begin{exercise}[$\star$]\label{ex:conpcomplete}
Show that if some NP-complete problem is in $\mathsf{coNP}$, then $\mathsf{NP}=\mathsf{coNP}$.
\end{exercise}

\begin{exercise}[$\star$]\label{ex:searchclique}
Suppose that an oracle answers, for any graph $G$ and number $k$, whether $G$ has a clique of size $k$. Show how to \emph{find} a largest clique of a graph with $n$ vertices using a polynomial number of questions.
\end{exercise}

\begin{exercise}[$\star$]\label{ex:hampath}
Show that deciding whether a graph has a Hamiltonian \emph{path} (visiting every vertex once, with arbitrary endpoints) is NP-complete.
\end{exercise}

\begin{exercise}[$\star$]\label{ex:kcol}
Show that $k$-colourability is NP-complete for every fixed $k\ge3$, and that $2$-colourability is in $\mathsf P$.
\end{exercise}

\begin{exercise}[$\star$]\label{ex:setcover}
Show that \textup{\textsc{Set cover}} (given subsets $S_1,\dots,S_r$ of a finite set $U$ and $k$, can $k$ of them cover $U$?) and \textup{\textsc{Dominating set}} (given $G$ and $k$, is there a set $D$ of $k$ vertices such that every vertex is in $D$ or adjacent to it?) are NP-complete.
\end{exercise}

\begin{exercise}[$\star\star$]\label{ex:partition}
Show that \textup{\textsc{Partition}} (can positive integers $a_1,\dots,a_N$ be split into two parts with equal sums?) is NP-complete.
\end{exercise}

\begin{exercise}[$\star\star$]\label{ex:binpacking}
In \emph{bin packing}, items of sizes $s_1,\dots,s_N\in(0,1]$ must be placed into as few bins of capacity $1$ as possible. Show that deciding whether two bins suffice is NP-complete, and deduce that unless $\mathsf P=\mathsf{NP}$ no polynomial algorithm always uses fewer than $\frac32$ times the optimal number of bins.
\end{exercise}

\begin{exercise}[$\star$]\label{ex:subgraphiso}
Show that \textup{\textsc{Subgraph isomorphism}} (given graphs $G$ and $H$, does $G$ contain a subgraph isomorphic to $H$?) is NP-complete.
\end{exercise}

\begin{exercise}[$\star$]\label{ex:longestpath}
Show that deciding whether a graph has a simple path of length at least $k$ between two given vertices is NP-complete. Shortest paths are found in polynomial time even with negative lengths on a graph without negative cycles (Tome~C2, Chapter~4); why does it not help to give every edge the length $-1$?
\end{exercise}

\begin{exercise}[$\star$]\label{ex:unarysubset}
Show that \textup{\textsc{Subset sum}} with the numbers written in unary is in $\mathsf P$. Why does this not contradict Theorem~\ref{thm:subsetsum}?
\end{exercise}

\begin{exercise}[$\star\star\star$]\label{ex:nexp}
Let $\mathsf{NEXP}=\bigcup_k\mathsf{NTIME}(2^{n^k})$. Show that if $\mathsf P=\mathsf{NP}$ then $\mathsf{EXP}=\mathsf{NEXP}$.
\end{exercise}

\section{Hints}

\begin{hint}{ex:conpcomplete}
Every problem in $\mathsf{NP}$ reduces to the complete one; take complements.
\end{hint}
\begin{hint}{ex:searchclique}
Find the size first; then try deleting vertices one at a time.
\end{hint}
\begin{hint}{ex:hampath}
Add a vertex to a graph, or split one into two copies, each with a new pendant vertex.
\end{hint}
\begin{hint}{ex:kcol}
Add $k-3$ new vertices joined to everything.
\end{hint}
\begin{hint}{ex:setcover}
For set cover, the sets are the stars of the vertices; for dominating set, subdivide each edge, keeping the old edges.
\end{hint}
\begin{hint}{ex:partition}
Add two numbers to a subset sum instance so that each part of a balanced partition must contain exactly one of them.
\end{hint}
\begin{hint}{ex:binpacking}
Scale a partition instance. An algorithm with ratio below $\frac32$ uses $2$ bins when $2$ suffice.
\end{hint}
\begin{hint}{ex:subgraphiso}
Take $H$ a complete graph.
\end{hint}
\begin{hint}{ex:longestpath}
Hamiltonian path between two vertices. With length $-1$ on every edge, cycles become negative.
\end{hint}
\begin{hint}{ex:unarysubset}
Dynamic programming over the achievable sums; compare the lengths.
\end{hint}
\begin{hint}{ex:nexp}
Pad: from $L\in\mathsf{NEXP}$ build $\{x\#1^{2^{|x|^k}}\}$.
\end{hint}

\section{Solutions}

\begin{solution}{ex:conpcomplete}
\idea Reductions respect complements, so one complete problem in $\mathsf{coNP}$ drags all of $\mathsf{NP}$ there.
\writeup Let $B$ be NP-complete and in $\mathsf{coNP}$. For $A\in\mathsf{NP}$, $A\le_pB$ by some $f$, and $\bar A\le_p\bar B$ by the same $f$. Since $\bar B\in\mathsf{NP}$, a certificate for $f(x)\in\bar B$ is a certificate for $x\in\bar A$ (Proposition~\ref{prop:certificates}, with the verifier first computing $f(x)$); so $\bar A\in\mathsf{NP}$ and $A\in\mathsf{coNP}$. Hence $\mathsf{NP}\subseteq\mathsf{coNP}$. Taking complements, $\mathsf{coNP}\subseteq\mathsf{NP}$, and the classes are equal.
\end{solution}

\begin{solution}{ex:searchclique}
\idea Ask for the size, then shrink the graph while the answer stays the same.
\writeup Find the largest $k$ such that $G$ has a clique of size $k$, with at most $n$ questions. Then for $v=v_1,\dots,v_n$ in turn: if the current graph without $v$ still has a clique of size $k$, delete $v$. At the end the remaining graph has a clique of size $k$, and every remaining vertex $v$ belongs to every such clique (deleting it destroyed them all when it was examined, and deleting other vertices later only removed cliques). So the remaining graph has exactly $k$ vertices, which form the clique. In all, at most $2n$ questions.
\end{solution}

\begin{solution}{ex:hampath}
\idea Turn a cycle into a path by cutting at a vertex, and force the path to end at the two halves of the cut.
\writeup The problem is in $\mathsf{NP}$. Reduce \textup{\textsc{Ham cycle}}: given $G$ and a vertex $v$, build $G'$ by adding a copy $v'$ of $v$, adjacent to the same vertices as $v$, and two new vertices $p$ and $p'$ adjacent only to $v$ and $v'$ respectively. A Hamiltonian cycle of $G$, cut at $v$, is $v,u_1,\dots,u_{n-1},v$; then $p,v,u_1,\dots,u_{n-1},v',p'$ is a Hamiltonian path of $G'$. Conversely, $p$ and $p'$ have degree $1$, so they are the endpoints of a Hamiltonian path of $G'$, which reads $p,v,u_1,\dots,u_{n-1},v',p'$; since $u_{n-1}$ is adjacent to $v'$, it is adjacent to $v$, and $v,u_1,\dots,u_{n-1},v$ is a Hamiltonian cycle of $G$ (for $n\ge3$).
\end{solution}

\begin{solution}{ex:kcol}
\idea Extra universal vertices use up extra colours.
\writeup Given $G$, add a clique $K$ of $k-3$ new vertices joined to every vertex of $G$. A $3$-colouring of $G$ extends with $k-3$ new colours. Conversely, in a $k$-colouring the vertices of $K$ have $k-3$ distinct colours that no vertex of $G$ may use, so $G$ is coloured with $3$ colours. For $2$-colourability: colour each component by breadth-first search, alternating colours by distance from a start vertex; the graph is $2$-colourable if and only if no edge joins two vertices of the same colour (an odd cycle otherwise appears), checked in linear time.
\end{solution}

\begin{solution}{ex:setcover}
\idea A vertex cover covers edges; a dominating set covers vertices, and subdividing edges turns one into the other.
\writeup Both problems are in $\mathsf{NP}$. Set cover: given $(G,k)$, let $U=E$ and, for each vertex $v$, $S_v$ the set of edges at $v$; $k$ sets cover $U$ if and only if $k$ vertices meet every edge. Dominating set: given $(G,k)$ with $G$ without isolated vertices, build $G'$ from $G$ by adding, for each edge $uv$, a new vertex $e_{uv}$ adjacent to $u$ and $v$. A vertex cover $C$ of $G$ dominates $G'$: every new vertex $e_{uv}$ has a neighbour in $C$, and every old vertex is in $C$ or adjacent to one (its edges are covered). Conversely, a dominating set $D$ of $G'$ can be moved off the new vertices: replace $e_{uv}\in D$ by $u$. The new set still dominates and has at most $|D|$ elements, and since each $e_{uv}$ is dominated by $u$ or $v$, it is a vertex cover of $G$.
\end{solution}

\begin{solution}{ex:partition}
\idea Two large numbers that cannot be in the same half fix which half carries the subset.
\writeup \textup{\textsc{Partition}} is in $\mathsf{NP}$. Given $a_1,\dots,a_N$ and $t$, with $S=\sum a_k$, add the numbers $b=2S-t$ and $c=S+t$. The total is $4S$, so a partition has two parts of sum $2S$. Since $b+c=3S>2S$, the two new numbers are in different parts, and the part containing $b$ is completed by old numbers of sum $2S-b=t$. Conversely, a subset of sum $t$, together with $b$, has sum $2S$, and the rest has sum $2S$ too. (If $t>S$, the instance is a no-instance and can be mapped to a fixed no-instance of \textup{\textsc{Partition}}.)
\end{solution}

\begin{solution}{ex:binpacking}
\idea Two bins of capacity $1$ for items of total size $2$ is exactly a partition.
\writeup The decision problem is in $\mathsf{NP}$. Given a partition instance with $S=\sum a_k$, take sizes $s_k=2a_k/S$, of total $2$ (if some $a_k>S/2$, there is no partition, and we output a fixed no-instance); two bins suffice if and only if the items split into two parts of size exactly $1$, that is, of sums $S/2$. Now suppose an algorithm always uses fewer than $\frac32\mathrm{OPT}$ bins. When $\mathrm{OPT}=2$ it uses fewer than $3$ bins, so exactly $2$; when $\mathrm{OPT}\ge3$ it uses at least $3$. So it decides partition, and $\mathsf P=\mathsf{NP}$.
\end{solution}

\begin{solution}{ex:subgraphiso}
\idea Clique is the special case where the pattern is complete.
\writeup The problem is in $\mathsf{NP}$: the certificate is an injective map from $H$ to $G$ preserving edges. The map $(G,k)\mapsto(G,K_k)$ reduces \textup{\textsc{Clique}} to it. By contrast, \emph{graph isomorphism}, where $H$ and $G$ have the same number of vertices and edges, is not known to be NP-complete (Section~\ref{sec:meaning}).
\end{solution}

\begin{solution}{ex:longestpath}
\idea A path of length $n-1$ is a Hamiltonian path; ``longest'' is not ``shortest with negated lengths'' because of cycles.
\writeup The problem is in $\mathsf{NP}$. Hamiltonian path between two given vertices $s$ and $t$ is NP-complete: in the solution of Exercise~\ref{ex:hampath}, the endpoints were forced to be $p$ and $p'$. A simple $s$--$t$ path of length at least $n-1$ in a graph with $n$ vertices visits every vertex. With lengths $-1$, a longest simple path is a shortest simple path, but every cycle becomes negative, and shortest path algorithms (Bellman and Ford) find shortest \emph{walks}, which are undefined in the presence of negative cycles; restricting to simple paths is the whole difficulty.
\end{solution}

\begin{solution}{ex:unarysubset}
\idea The dynamic programme is polynomial in the values, and unary makes values as long as the input.
\writeup Let $R_k$ be the set of sums of subsets of $a_1,\dots,a_k$ that are at most $t$: $R_0=\{0\}$ and $R_k=R_{k-1}\cup\{r+a_k:r\in R_{k-1},\ r+a_k\le t\}$. Each $R_k$ has at most $t+1$ elements, so the whole computation takes time $O(Nt)$, and the answer is whether $t\in R_N$. In unary, the input has length at least $t+N$, so $O(Nt)$ is polynomial in it. The reduction of Theorem~\ref{thm:subsetsum} produces numbers with $n+m$ decimal digits; written in unary, they would have length about $10^{n+m}$, so it is not a polynomial reduction to the unary problem. Subset sum is NP-complete, but not strongly.
\end{solution}

\begin{solution}{ex:nexp}
\idea Padding stretches an exponential-time problem into a polynomial-time one, where the hypothesis applies.
\writeup Let $L\in\mathsf{NTIME}(2^{n^k})$, decided by a nondeterministic machine $N$. Let $L'=\{x\#1^{2^{|x|^k}}:x\in L\}$. On an input of length $m$, check the format and run $N$ on $x$; this takes nondeterministic time $O(2^{|x|^k})=O(m)$, so $L'\in\mathsf{NP}=\mathsf P$, decided by a deterministic machine $M$ in time $m^c$. To decide $L$ deterministically, on input $x$ write $x\#1^{2^{|x|^k}}$, in time $O(2^{|x|^k})$, and run $M$, in time $O(2^{c|x|^k})$. So $L\in\mathsf{EXP}$, and $\mathsf{NEXP}\subseteq\mathsf{EXP}$; the other inclusion is clear. Read backwards: if $\mathsf{EXP}\ne\mathsf{NEXP}$, then $\mathsf P\ne\mathsf{NP}$. Equalities propagate upwards, separations downwards.
\end{solution}

\section*{Sources and further reading}

\begin{itemize}
\item \emph{A research paper to read.} R.~M.~Karp, \emph{Reducibility among combinatorial problems}, in R.~E.~Miller and J.~W.~Thatcher (eds), \emph{Complexity of Computer Computations}, Plenum Press, 1972, 85--103. The paper that showed how far NP-completeness reaches. The first sections define the classes $\mathcal P$ and $\mathcal{NP}$ and the reducibility we call $\le_p$, and recall Cook's theorem. Then comes the main theorem, a list of twenty-one problems with a tree of reductions much like Figure~\ref{fig:reductions}: satisfiability, \textup{\textsc{0-1 integer programming}}, \textup{\textsc{Clique}}, \textup{\textsc{Set packing}}, \textup{\textsc{Node cover}}, \textup{\textsc{Set covering}}, \textup{\textsc{Feedback node set}}, \textup{\textsc{Feedback arc set}}, \textup{\textsc{Directed}} and \textup{\textsc{Undirected Hamilton circuit}}, \textup{\textsc{Satisfiability with at most 3 literals per clause}}, \textup{\textsc{Chromatic number}}, \textup{\textsc{Clique cover}}, \textup{\textsc{Exact cover}}, \textup{\textsc{Hitting set}}, \textup{\textsc{Steiner tree}}, \textup{\textsc{3-dimensional matching}}, \textup{\textsc{Knapsack}}, \textup{\textsc{Job sequencing}}, \textup{\textsc{Partition}}, \textup{\textsc{Max cut}}. Each reduction takes a few lines, sometimes only a formula; checking them, with the proofs of this chapter as a model, is an excellent exercise. Karp ends with three problems he could not classify: graph isomorphism, the recognition of composite numbers, and linear inequalities. Linear programming was put in $\mathsf P$ by Khachiyan in 1979 (Tome~C2, Chapter~7), primality in 2002; graph isomorphism is still unclassified.
\item S.~A.~Cook, \emph{The complexity of theorem-proving procedures}, Proceedings of the Third Annual ACM Symposium on Theory of Computing (1971), 151--158. Theorem~1 is our Theorem~\ref{thm:cooklevin}, in its form for tautologies; read it to see the tableau in its original notation. L.~A.~Levin, \emph{Universal sequential search problems}, Problems of Information Transmission 9 (1973), 265--266 (translated from the Russian).
\item \emph{The stories.} J.~Hartmanis, \emph{Gödel, von Neumann and the P=?NP problem}, Bulletin of the EATCS 38 (1989), 101--107, which publishes Gödel's letter; M.~Sipser, \emph{The history and status of the P versus NP question}, Proceedings of the 24th ACM Symposium on Theory of Computing (1992), 603--618, with a translation of the letter; R.~M.~Karp, \emph{Combinatorics, complexity, and randomness}, Communications of the ACM 29 (1986), 98--109, his Turing Award lecture, on Cook and Berkeley; S.~A.~Cook, \emph{The P versus NP problem}, the official problem description for the Clay Mathematics Institute (2000).
\item M.~R.~Garey and D.~S.~Johnson, \emph{Computers and Intractability: A Guide to the Theory of NP-Completeness}, Freeman, 1979, with its catalogue of problems; M.~R.~Garey and D.~S.~Johnson, \emph{``Strong'' NP-completeness results: motivation, examples, and implications}, Journal of the ACM 25 (1978), 499--508; M.~R.~Garey and D.~S.~Johnson, \emph{Complexity results for multiprocessor scheduling under resource constraints}, SIAM Journal on Computing 4 (1975), 397--411.
\item L.~Stockmeyer, \emph{Planar 3-colorability is polynomial complete}, SIGACT News 5 (1973), 19--25; C.~H.~Papadimitriou, \emph{On the complexity of integer programming}, Journal of the ACM 28 (1981), 765--768; B.~Mohar, \emph{Isoperimetric numbers of graphs}, Journal of Combinatorial Theory, Series B 47 (1989), 274--291; R.~E.~Ladner, \emph{On the structure of polynomial time reducibility}, Journal of the ACM 22 (1975), 155--171; L.~Babai, \emph{Graph isomorphism in quasipolynomial time}, Proceedings of the 48th ACM Symposium on Theory of Computing (2016), 684--697.
\item M.~Sipser, \emph{Introduction to the Theory of Computation}, 3rd edition, Cengage, 2013, Chapter~7; S.~Arora and B.~Barak, \emph{Computational Complexity: A Modern Approach}, Cambridge University Press, 2009, Chapter~2.
\end{itemize}
